\chapter[King Louis XIV at the Age of Twenty-Two]{What His Majesty King Louis XIV Was at the Age of Twenty-Two}
\chaptermark{King Louis XIV at the Age of Twenty-Two}

	\lettrine{I}{t} has been seen, by the account we have endeavoured to  give of it, that the entrée of King Louis XIV into the city of Blois had been  noisy and brilliant; his young majesty had therefore appeared perfectly  satisfied with it.

\zz
On arriving beneath the porch of the Castle of the States, the king met,  surrounded by his guards and gentlemen, with S. A. R. the duke, Gaston of  Orleans, whose physiognomy, naturally rather majestic, had borrowed on this  solemn occasion a fresh lustre and a fresh dignity. On her part, Madame,  dressed in her robes of ceremony, awaited, in the interior balcony, the  entrance of her nephew. All the windows of the old castle, so deserted and  dismal on ordinary days, were resplendent with ladies and lights.

It was then to the sound of drums, trumpets, and vivats, that the young king  crossed the threshold of that castle in which, seventy-two years before, Henry  III had called in the aid of assassination and treachery to keep upon his head  and in his house a crown which was already slipping from his brow, to fall into  another family.

All eyes, after having admired the young king, so handsome and so agreeable,  sought for that other king of France, much otherwise king than the former, and  so old, so pale, so bent, that people called the Cardinal Mazarin.

Louis was at this time endowed with all the natural gifts which make the  perfect gentleman; his eye was brilliant, mild, and of a clear azure blue. But  the most skilful physiognomists, those divers into the soul, on fixing their  looks upon it, if it had been possible for a subject to sustain the glance of  the king,—the most skilful physiognomists, we say, would never have been  able to fathom the depths of that abyss of mildness. It was with the eyes of  the king as with the immense depths of the azure heavens, or with those more  terrific, and almost as sublime, which the Mediterranean reveals under the  keels of its ships in a clear summer day, a gigantic mirror in which heaven  delights to reflect sometimes its stars, sometimes its storms.

The king was short of stature—he was scarcely five feet two inches: but  his youth made up for this defect, set off likewise by great nobleness in all  his movements, and by considerable address in all bodily exercises.

Certes, he was already quite a king, and it was a great thing to be a king in  that period of traditional devotedness and respect; but as, up to that time, he  had been but seldom and always poorly shown to the people, as they to whom he  was shown saw him by the side of his mother, a tall woman, and monsieur le  cardinal, a man of commanding presence, many found him so little of a king as  to say,—

<Why, the king is not so tall as monsieur le cardinal!>

Whatever may be thought of these physical observations, which were principally  made in the capital, the young king was welcomed as a god by the inhabitants of  Blois, and almost like a king by his uncle and aunt, Monsieur and Madame, the  inhabitants of the castle.

It must, however, be allowed, that when he saw, in the hall of reception,  chairs of equal height for himself, his mother, the cardinal, and his uncle and  aunt, a disposition artfully concealed by the semi-circular form of the  assembly, Louis XIV became red with anger, and looked around him to ascertain  by the countenances of those that were present, if this humiliation had been  prepared for him. But as he saw nothing upon the impassible visage of the  cardinal, nothing on that of his mother, nothing on those of the assembly, he  resigned himself, and sat down, taking care to be seated before anybody else.

The gentlemen and ladies were presented to their majesties and monsieur le  cardinal.

The king remarked that his mother and he scarcely knew the names of any of the  persons who were presented to them; whilst the cardinal, on the contrary, never  failed, with an admirable memory and presence of mind, to talk to every one  about his estates, his ancestors, or his children, some of whom he named, which  enchanted those worthy country gentlemen, and confirmed them in the idea that  he alone is truly king who knows his subjects, from the same reason that the  sun has no rival, because the sun alone warms and lightens.

The study of the young king, which had begun a long time before, without  anybody suspecting it, was continued then, and he looked around him attentively  to endeavour to make out something in the physiognomies which had at first  appeared the most insignificant and trivial.

A collation was served. The king, without daring to call upon the hospitality  of his uncle, had waited for it impatiently. This time, therefore, he had all  the honours due, if not to his rank, at least to his appetite.

As to the cardinal, he contented himself with touching with his withered lips a  bouillon, served in a golden cup. The all-powerful minister, who had taken her  regency from the queen, and his royalty from the king, had not been able to  take a good stomach from nature.

Anne of Austria, already suffering from the cancer which six or eight years  after caused her death, ate very little more than the cardinal.

For Monsieur, already puffed up with the great event which had taken place in  his provincial life, he ate nothing whatever.

Madame alone, like a true Lorrainer, kept pace with his majesty; so that Louis  XIV, who, without this partner, might have eaten nearly alone, was at first  much pleased with his aunt, and afterwards with M. de Saint-Remy, her maître d'hôtel, who had really distinguished himself.

The collation over, at a sign of approbation from M. de Mazarin, the king  arose, and, at the invitation of his aunt, walked about among the ranks of the  assembly.

The ladies then observed—there are certain things for which women are as  good observers at Blois as at Paris—the ladies then observed that Louis  XIV had a prompt and bold look, which premised a distinguished appreciator of  beauty. The men, on their part, observed that the prince was proud and haughty,  that he loved to look down those who fixed their eyes upon him too long or too  earnestly, which gave presage of a master.

Louis XIV had accomplished about a third of his review when his ears were  struck with a word which his eminence pronounced whilst conversing with  Monsieur.

This word was the name of a woman.

Scarcely had Louis XIV heard this word than he heard, or rather listening to  nothing else; and neglecting the arc of the circle which awaited his visit, his  object seemed to be to come as quickly as possible to the extremity of the  curve.

Monsieur, like a good courtier, was inquiring of monsieur le cardinal after the  health of his nieces; he regretted, he said, not having the pleasure of  receiving them at the same time with their uncle; they must certainly have  grown in stature, beauty and grace, as they had promised to do the last time  Monsieur had seen them.

What had first struck the king was a certain constraint in the voices of the  two interlocutors. The voice of Monsieur was calm and natural when he spoke  thus; while that of M. de Mazarin jumped by a note and a half to reply above  the diapason of his usual voice. It might have been said that he wished that  voice to strike, at the end of the salon, any ear that was too distant.

<Monseigneur,> replied he, <Mesdemoiselles de Mazarin have  still to finish their education: they have duties to fulfil, and a position to  make. An abode in a young and brilliant court would dissipate them a  little.>

Louis, at this last sentence, smiled sadly. The court was young, it was true,  but the avarice of the cardinal had taken good care that it should not be  brilliant.

<You have nevertheless no intention,> replied Monsieur, <to  cloister them or make them bourgeoises?>

<Not at all,> replied the cardinal, forcing his Italian  pronunciation in such a manner that, from soft and velvety as it was, it became  sharp and vibrating; <not at all: I have a full and fixed intention to  marry them, and that as well as I shall be able.>

<Parties will not be wanting, monsieur le cardinal,> replied  Monsieur, with a bonhomie worthy of one tradesman congratulating another.

<I hope not, monseigneur, and with reason, as God has been pleased to  give them grace, intelligence, and beauty.>

During this conversation, Louis XIV, conducted by Madame, accomplished, as we  have described, the circle of presentations.

<Mademoiselle Auricule,> said the princess, presenting to his  majesty a fat, fair girl of two-and-twenty, who at a village fete might have  been taken for a peasant in Sunday finery,—<the daughter of my  music-mistress.>

The king smiled. Madame had never been able to extract four correct notes from  either viol or harpsichord.

<Mademoiselle Aure de Montalais,> continued Madame; <a young  lady of rank, and my good attendant.>

This time it was not the king that smiled; it was the young lady presented,  because, for the first time in her life, she heard, given to her by Madame, who  generally showed no tendency to spoil her, such an honourable qualification.

Our old acquaintance Montalais, therefore, made his majesty a profound  courtesy, the more respectful from the necessity she was under of concealing  certain contractions of her laughing lips, which the king might not have  attributed to their real cause.

It was just at this moment that the king caught the word which startled him.

<And the name of the third?> asked Monsieur.

<Mary, monseigneur,> replied the cardinal.

There was doubtless some magical influence in that word, for, as we have said,  the king started in hearing it, and drew Madame towards the middle of the  circle, as if he wished to put some confidential question to her, but, in  reality, for the sake of getting nearer to the cardinal.

<Madame, my aunt,> said he, laughing, and in a suppressed voice,  <my geography-master did not teach me that Blois was at such an immense  distance from Paris.>

<What do you mean, nephew?> asked Madame.

<Why, because it would appear that it requires several years, as regards  fashion, to travel the distance!—Look at those young ladies!>

<Well; I know them all.>

<Some of them are pretty.>

<Don't say that too loud, monsieur my nephew; you will drive them  wild.>

<Stop a bit, stop a bit, dear aunt!> said the king, smiling;  <for the second part of my sentence will serve as a corrective to the  first. Well, my dear aunt, some of them appear old and others ugly, thanks to  their ten-year-old fashions.>

<But, sire, Blois is only five days' journey from Paris.>

<Yes, that is it,> said the king: <two years behind for each  day.>

<Indeed! do you really think so? Well, that is strange! It never struck  me.>

<Now, look, aunt,> said Louis XIV, drawing still nearer to  Mazarin, under the pretext of gaining a better point of view, <look at  that simple white dress by the side of those antiquated specimens of finery,  and those pretentious coiffures. She is probably one of my mother's maids  of honour, though I don't know her.>

<Ah! ah! my dear nephew!> replied Madame, laughing; <permit  me to tell you that your divinatory science is at fault for once. The young  lady you honour with your praise is not a Parisian, but a Blaisoise.>

<Oh, aunt!> replied the king with a look of doubt.

<Come here, Louise,> said Madame.

And the fair girl, already known to you under that name, approached them,  timid, blushing, and almost bent beneath the royal glance.

<Mademoiselle Louise Françoise de la Beaume le Blanc, the daughter of the  Marquise de la Vallière,> said Madame, ceremoniously.

The young girl bowed with so much grace, mingled with the profound timidity  inspired by the presence of the king, that the latter lost, while looking at  her, a few words of the conversation of Monsieur and the cardinal.

<Daughter-in-law,> continued Madame, <of M. de Saint-Remy, my  maître d'hôtel, who presided over the confection of that excellent daube  truffée which your majesty seemed so much to appreciate.>

No grace, no youth, no beauty, could stand out against such a presentation. The  king smiled. Whether the words of Madame were a pleasantry, or uttered in all  innocency, they proved the pitiless immolation of everything that Louis had  found charming or poetic in the young girl. Mademoiselle de la Vallière, for  Madame and, by rebound, for the king, was, for a moment, no more than the  daughter of a man of a superior talent over dindes truffées.

But princes are thus constituted. The gods, too, were just like this in  Olympus. Diana and Venus, no doubt, abused the beautiful Alcmena and poor Io,  when they condescended for distraction's sake, to speak, amidst nectar  and ambrosia, of mortal beauties, at the table of Jupiter.

Fortunately, Louise was so bent in her reverential salute, that she did not  catch either Madame's words or the king's smile. In fact, if the  poor child, who had so much good taste as alone to have chosen to dress herself  in white amidst all her companions—if that dove's heart, so easily  accessible to painful emotions, had been touched by the cruel words of Madame,  or the egotistical cold smile of the king, it would have annihilated her.

And Montalais herself, the girl of ingenious ideas, would not have attempted to  recall her to life; for ridicule kills beauty even.

But fortunately, as we have said, Louise, whose ears were buzzing, and her eyes  veiled by timidity,—Louise saw nothing and heard nothing; and the king,  who had still his attention directed to the conversation of the cardinal and  his uncle, hastened to return to them.

He came up just at the moment Mazarin terminated by saying: <Mary, as  well as her sisters, has just set off for Brouage. I make them follow the  opposite bank of the Loire to that along which we have travelled; and if I  calculate their progress correctly, according to the orders I have given, they  will to-morrow be opposite Blois.>

These words were pronounced with that tact—that measure, that  distinctness of tone, of intention, and reach—which made del Signor  Giulio Mazarini the first comedian in the world.

It resulted that they went straight to the heart of Louis XIV, and the  cardinal, on turning round at the simple noise of the approaching footsteps of  his majesty, saw the immediate effect of them upon the countenance of his  pupil, an effect betrayed to the keen eyes of his eminence by a slight increase  of colour. But what was the ventilation of such a secret to him whose craft had  for twenty years deceived all the diplomatists of Europe?

From the moment the young king heard these last words, he appeared as if he had  received a poisoned arrow in his heart. He could not remain quiet in a place,  but cast around an uncertain, dead, and aimless look over the assembly. He with  his eyes interrogated his mother more than twenty times: but she, given up to  the pleasure of conversing with her sister-in-law, and likewise constrained by  the glance of Mazarin, did not appear to comprehend any of the supplications  conveyed by the looks of her son.

From this moment, music, lights, flowers, beauties, all became odious and  insipid to Louis XIV After he had a hundred times bitten his lips, stretched  his legs and his arms like a well-brought-up child, who, without daring to  gape, exhausts all the modes of evincing his weariness—after having  uselessly again implored his mother and the minister, he turned a despairing  look towards the door, that is to say, towards liberty.

At this door, in the embrasure of which he was leaning, he saw, standing out  strongly, a figure with a brown and lofty countenance, an aquiline nose, a  stern but brilliant eye, gray and long hair, a black moustache, the true type of  military beauty, whose gorget, more sparkling than a mirror, broke all the  reflected lights which concentrated upon it, and sent them back as lightning.  This officer wore his gray hat with its long red plumes upon his head, a proof  that he was called there by his duty, and not by his pleasure. If he had been  brought thither by his pleasure—if he had been a courtier instead of a  soldier, as pleasure must always be paid for at the same price—he would  have held his hat in his hand.

That which proved still better that this officer was upon duty, and was  accomplishing a task to which he was accustomed, was, that he watched, with  folded arms, remarkable indifference, and supreme apathy, the joys and ennuis  of this fete. Above all, he appeared, like a philosopher, and all old soldiers  are philosophers,—he appeared above all to comprehend the ennuis  infinitely better than the joys; but in the one he took his part, knowing very  well how to do without the other.

Now, he was leaning, as we have said, against the carved door-frame when the  melancholy, weary eyes of the king, by chance, met his.

It was not the first time, as it appeared, that the eyes of the officer had met  those eyes, and he was perfectly acquainted with the expression of them; for,  as soon as he had cast his own look upon the countenance of Louis XIV, and had  read by it what was passing in his heart—that is to say, all the ennui  that oppressed him—all the timid desire to go out which agitated  him,—he perceived he must render the king a service without his  commanding it,—almost in spite of himself. Boldly, therefore, as if he  had given the word of command to cavalry in battle, <On the king's  service!> cried he, in a clear, sonorous voice.

At these words, which produced the effect of a peal of thunder, prevailing over  the orchestra, the singing and the buzz of the promenaders, the cardinal and  the queen-mother looked at each other with surprise.

Louis XIV, pale, but resolved, supported as he was by that intuition of his  own thought which he had found in the mind of the officer of musketeers, and  which he had just manifested by the order given, arose from his chair, and took  a step towards the door.

<Are you going, my son?> said the queen, whilst Mazarin satisfied  himself with interrogating by a look which might have appeared mild if it had  not been so piercing.

<Yes, madame,> replied the king; <I am fatigued, and,  besides, wish to write this evening.>

A smile stole over the lips of the minister, who appeared, by a bend of the  head, to give the king permission.

Monsieur and Madame hastened to give orders to the officers who presented  themselves.

The king bowed, crossed the hall, and gained the door, where a hedge of twenty  musketeers awaited him. At the extremity of this hedge stood the officer,  impassible, with his drawn sword in his hand. The king passed, and all the  crowd stood on tip-toe, to have one more look at him.

Ten musketeers, opening the crowd of the ante-chambers and the steps, made way  for his majesty. The other ten surrounded the king and Monsieur, who had  insisted upon accompanying his majesty. The domestics walked behind. This  little cortège escorted the king to the chamber destined for him. The apartment  was the same that had been occupied by Henry III during his sojourn in the  States.

Monsieur had given his orders. The musketeers, led by their officer, took  possession of the little passage by which one wing of the castle communicates  with the other. This passage was commenced by a small square ante-chamber, dark  even in the finest days. Monsieur stopped Louis XIV.

<You are passing now, sire,> said he, <the very spot where  the Duc de Guise received the first stab of the poniard.>

The king was ignorant of all historical matters; he had heard of the fact, but  he knew nothing of the localities or the details.

<Ah!> said he with a shudder.

And he stopped. The rest, both behind and before him, stopped likewise.

<The duc, sire,> continued Gaston, <was nearly were I stand:  he was walking in the same direction as your majesty; M. de Loignac was exactly  where your lieutenant of musketeers is; M. de Saint-Maline and his  majesty's ordinaries were behind him and around him. It was here that he  was struck.>

The king turned towards his officer, and saw something like a cloud pass over  his martial and daring countenance.

<Yes, from behind!> murmured the lieutenant, with a gesture of  supreme disdain. And he endeavoured to resume the march, as if ill at ease at  being between walls formerly defiled by treachery.

But the king, who appeared to wish to be informed, was disposed to give another  look at this dismal spot.

Gaston perceived his nephew's desire.

<Look, sire,> said he, taking a flambeau from the hands of M. de  Saint-Remy, <this is where he fell. There was a bed there, the curtains  of which he tore with catching at them.>

<Why does the floor seem hollowed out at this spot?> asked Louis.

<Because it was here the blood flowed,> replied Gaston; <the  blood penetrated deeply into the oak, and it was only by cutting it out that  they succeeded in making it disappear. And even then,> added Gaston,  pointing the flambeau to the spot, <even then this red stain resisted  all the attempts made to destroy it.>

Louis XIV raised his head. Perhaps he was thinking of that bloody trace that  had once been shown him at the Louvre, and which, as a pendant to that of  Blois, had been made there one day by the king his father with the blood of  Concini.

<Let us go on,> said he.

The march was resumed promptly; for emotion, no doubt, had given to the voice  of the young prince a tone of command which was not customary with him. When he  arrived at the apartment destined for the king, which communicated not only  with the little passage we have passed through, but further with the great  staircase leading to the court,—

<Will your majesty,> said Gaston, <condescend to occupy this  apartment, all unworthy as it is to receive you?>

<Uncle,> replied the young king, <I render you my thanks for  your cordial hospitality.>

Gaston bowed to his nephew, embraced him, and then went out.

Of the twenty musketeers who had accompanied the king, ten reconducted Monsieur  to the reception-rooms, which were not yet empty, notwithstanding the king had  retired.

The ten others were posted by their officer, who himself explored, in five  minutes, all the localities, with that cold and certain glance which not even  habit gives unless that glance belongs to genius.

Then, when all were placed, he chose as his headquarters the ante-chamber, in  which he found a large fauteuil, a lamp, some wine, some water, and some dry  bread.

He refreshed his lamp, drank half a glass of wine, curled his lip with a smile  full of expression, installed himself in his large armchair, and made  preparations for sleeping.  